%! Author = lili
%! Date = 4/27/26
\section{Präsenzaufgaben 1.1.2}
\subsection{Aufgabe 1.2.4}

Die Funktion $f: \rz \rightarrow \rz$ sei gegeben durch
\[
    f(x) = \begin{cases}
   3cx^2(1-x) & \text{für } x \in [0,1],
    \\ 0 & \text{sonst}
\end{cases}
\]
wobei $c \in \rz$ eine Konstante sei.
\begin{auf}
    \subsubsection*{(a) Bestimmen Sie ein $c \in \rz$ derart, dass $f$ eine Dichte ist}
    \paragraph{Lösung:} $f(x)$ ist eine Dichte, wenn $\int_{\rz}  f(x) \dx = 1$ also $\int_{-\infty}^{\infty}  f(x) \dx = 1$
% TODO tag einbauen?
    (Von Studi vorgestellt)
    \begin{align*}
        \int 3cx^2(1-x) \dx & = 3c \cdot \int x^2(1-x) \dx \\
        & = 3c \cdot \int_0^1 x^2-x^3 \dx \\
        & = 3c \cdot \Big[ \frac{1}{3} x^3  - \frac{1}{4} x^4 \Big]^1_0 \\
        & = 3c \cdot (\frac{1}{3} \cdot \frac{1}{4} \\
        & = c \cdot \frac{1}{4} \\ \\
        c\cdot \frac{1}{4} = 1 & \Rightarrow c = 4
    \end{align*}
\end{auf}
Sei $\pe$ das Wahrscheinlichkeitsmaß mit Dichte $f$, wie Sie sie in (a) bestimmt haben. \\
\textit{Hinweis:} Falls Sie die Konstante c nicht bestimmen konnten, rechnen Sie im folgenden bitte
mit der unbestimmten Konstanten weiter.

\begin{auf}
    \subsubsection*{(b) Berechnen Sie die folgenden Wahrscheinlichkeiten}
    \[
        \pe((\frac{1}{2}, \infty)), \pe((- \infty, 0)) , \pe([1, \infty)) \text{ und } \pe(\Big\{\frac{1}{2}\Big\})
    \]
    \paragraph{Lösung:}
    (von Tutor)
    \begin{align*}
        pe((\frac{1}{2}, \infty)) & = \frac{11}{16} \\ \\
        \pe((- \infty, 0)) & = \int^0_{-\infty} f(x) \dx \\
        & = \int^0_{-\infty} 0 \dx \\ & = 0 \\ \\
        \pe([1, \infty)) & = \int^{\infty}_1 f(x) \dx \\
        & = \int^{\infty}_1 0 \dx \\
        & = 0 \\ \\
        \pe(\Big\{\frac{1}{2}\Big\}) & = \int^{\frac{1}{2}}_{\frac{1}{2}} f(x) \dx \\
        &= 12 \cdot \Big[\frac{1}{3} \cdot x^3 - \frac{1}{4} \cdot x^4 \Big]^{\frac{1}{2}}_{\frac{1}{2}} \\
        & = 12 \cdot (\frac{1}{3} \cdot \frac{1}{2^3} - \frac{1}{4} \cdot \frac{1}{2^4} - \frac{1}{3} \cdot
        \frac{1}{2^3} + \frac{1}{4} \cdot \frac{1}{2^4}) \\ % TODO Rechnung kontrollieren wegen +/-
        & = 0
    \end{align*}
\end{auf}
